\documentclass[11pt,a4paper]{article}

\usepackage[margin=2.5cm]{geometry}
\usepackage[T1]{fontenc}
\usepackage[utf8]{inputenc}
\usepackage{lmodern}
\usepackage{amsmath,amssymb,bm}
\usepackage{booktabs}
\usepackage{array}
\usepackage{xcolor}
\usepackage{hyperref}
\usepackage{enumitem}
\usepackage{listings}
\usepackage{longtable}

\hypersetup{
  colorlinks=true,
  linkcolor=blue!55!black,
  urlcolor=blue!55!black,
  citecolor=blue!55!black
}

\lstdefinestyle{shellstyle}{
  basicstyle=\ttfamily\small,
  breaklines=true,
  frame=single,
  columns=fullflexible,
  backgroundcolor=\color{gray!6}
}
\lstset{style=shellstyle}

\newcommand{\ket}[1]{\left|#1\right\rangle}
\newcommand{\bra}[1]{\left\langle #1\right|}
\newcommand{\braket}[2]{\left\langle #1\middle|#2\right\rangle}
\newcommand{\dd}{\mathrm{d}}
\newcommand{\ii}{\mathrm{i}}
\newcommand{\Tr}{\mathrm{Tr}}

\title{A/B Interface Chain Model\\
\large Theoretical Background, Hamiltonian Definition, Time Propagation, and User Guide}
\author{AB Interface Toy-Model Software Documentation}
\date{Version for the no-impact GUI package}

\begin{document}
\maketitle

\tableofcontents
\newpage

\section{Purpose of the Software}

This software implements a compact density-matrix toy model for short-time laser-driven carrier and spin redistribution at a non-matching interface between two materials, denoted by $A$ and $B$.  The model is designed for qualitative mechanism analysis rather than quantitative materials prediction.  It is useful when one wants to test whether a proposed sequence of optical excitation, spin-orbit-induced spin mixing, interlayer transfer, and intramaterial relaxation can generate a target population-flow pattern.

A useful effective interpretation targeted by the default parameter regime is the emergent chain-like redistribution pattern
\begin{equation}
A_v\uparrow
\;\xrightarrow{\text{laser}}\;
A_c\uparrow
\;\xrightarrow{\text{SOC mixing in }A_c}\;
A_c\downarrow
\;\xrightarrow{\text{spin-conserving interlayer transfer}}\;
B_c\downarrow
\;\xrightarrow{\text{intramaterial relaxation in }B}\;
B_v\downarrow.
\end{equation}
This sequence is \emph{not} hard-coded as a single ordered path.  The code instead implements elementary optical, SOC-mixing, interlayer-transfer, and intramaterial-relaxation channels.  Under suitable energy alignment and rate parameters, the occupation dynamics may be interpreted as the above chain-like redistribution.  This replaces the earlier, less transparent direct interface-SOC transition
\begin{equation}
A_c\uparrow \leftrightarrow B_v\downarrow,
\end{equation}
which bundled material change, spin flip, and band relaxation into a single phenomenological matrix element.  In the present implementation, spin mixing is produced by material-internal SOC terms, while the interlayer transfer is spin-conserving.

The code generates time-dependent projected occupations, magnetic moment proxies, static-eigenbasis occupations, state-resolved occupation-change diagrams, CSV data, and a graphical user interface.  It is packaged as a Python project and can be run either from the command line or through a Tkinter desktop GUI.

\section{Physical Model Overview}

The model separates the dynamics into coherent Hamiltonian evolution and incoherent post-pulse rate processes.  The time-dependent density matrix is denoted by $\rho(t)$.  The coherent part is governed by
\begin{equation}
H(t) = H_{\mathrm{static}} + H_{\mathrm{laser}}(t),
\end{equation}
where
\begin{equation}
H_{\mathrm{static}}
=
H_{\mathrm{bands}}
+
H_{\mathrm{SOC,mix}}
+
H_{\mathrm{interlayer}}.
\end{equation}

The model contains the following physical components:
\begin{enumerate}[leftmargin=2em]
  \item \textbf{No-SOC non-matching bands}: materials $A$ and $B$ have independent gaps, energy offsets, and band widths.
  \item \textbf{SOC-derived diagonal splitting}: spin splitting is not an independent input; it is computed from material-internal SOC strengths $\lambda_{\mathrm{soc},A}$ and $\lambda_{\mathrm{soc},B}$.
  \item \textbf{Material-internal SOC mixing}: off-diagonal couplings mix $\uparrow$ and $\downarrow$ components within the same material and band manifold.
  \item \textbf{Spin-conserving interlayer hopping}: the primary interlayer process is $A_c s \leftrightarrow B_c s$ with $s=\uparrow,\downarrow$.
  \item \textbf{Laser excitation}: the optical pulse drives spin-conserving intramaterial $v\leftrightarrow c$ transitions.
  \item \textbf{Post-pulse rate processes}: after the laser pulse, the model applies spin-conserving interlayer downhill transfer and intramaterial $c\rightarrow v$ relaxation.
\end{enumerate}

\subsection*{Important Modeling Clarification}

The software does not propagate classical single-electron trajectories and does not prescribe that an electron must follow one specific sequence.  It propagates a density matrix and applies a set of allowed elementary couplings and rate channels.  Therefore statements such as
\begin{equation}
A_v\uparrow\rightarrow A_c\uparrow\rightarrow A_c\downarrow\rightarrow B_c\downarrow\rightarrow B_v\downarrow
\end{equation}
should be read as an \emph{effective pathway interpretation} of projected occupation changes, not as a hard-coded algorithmic transition chain.  Directly computed observables are occupations, coherences, and their projections onto material, spin, band, and static-eigenstate subspaces.

The former impact-excitation branch has been removed from the main model.  Without explicit electron-electron Coulomb matrix elements and a controlled bath spectral density, impact excitation is too phenomenological for the present short-time emergent chain-like redistribution scenario and obscures pathway interpretation.

\section{Single-Particle Basis}

Each material and spin block contains $N$ single-particle states.  The total Hilbert-space dimension is
\begin{equation}
D = 4N.
\end{equation}
The basis ordering is
\begin{equation}
A_\uparrow,
\quad
A_\downarrow,
\quad
B_\uparrow,
\quad
B_\downarrow.
\end{equation}
In each block, the first $N/2$ states are valence-like and the last $N/2$ states are conduction-like.  Therefore $N$ must be even.  In index form,
\begin{center}
\begin{tabular}{lll}
\toprule
Block & Index range & Meaning \\
\midrule
$A_\uparrow$ & $0,\ldots,N-1$ & material $A$, spin up \\
$A_\downarrow$ & $N,\ldots,2N-1$ & material $A$, spin down \\
$B_\uparrow$ & $2N,\ldots,3N-1$ & material $B$, spin up \\
$B_\downarrow$ & $3N,\ldots,4N-1$ & material $B$, spin down \\
\bottomrule
\end{tabular}
\end{center}
Within any block,
\begin{equation}
0,\ldots,N/2-1 \quad \text{are valence-like states},
\end{equation}
while
\begin{equation}
N/2,\ldots,N-1 \quad \text{are conduction-like states}.
\end{equation}
The software uses helper functions in \texttt{basis.py} to avoid hard-coded manual indexing.

\section{No-SOC Band Energies}

The no-SOC band structure is specified by material-dependent gaps, offsets, and bandwidths.  For material $X\in\{A,B\}$, the valence-band maximum and conduction-band minimum before SOC splitting are approximately
\begin{equation}
E^{(0)}_{X,\mathrm{VBM}} = \mathrm{offset}_X - \frac{1}{2}\mathrm{gap}_X,
\end{equation}
\begin{equation}
E^{(0)}_{X,\mathrm{CBM}} = \mathrm{offset}_X + \frac{1}{2}\mathrm{gap}_X.
\end{equation}
The valence ladder is constructed downward from the VBM:
\begin{equation}
E^{(0)}_{X,v,m}=E^{(0)}_{X,\mathrm{VBM}}-\epsilon^{X,v}_m,
\end{equation}
where $\epsilon^{X,v}_m$ spans the valence bandwidth.  The conduction ladder is constructed upward from the CBM:
\begin{equation}
E^{(0)}_{X,c,n}=E^{(0)}_{X,\mathrm{CBM}}+\epsilon^{X,c}_n,
\end{equation}
where $\epsilon^{X,c}_n$ spans the conduction bandwidth.

The default alignment is staggered so that material $A$ has the lower conduction-band edge and material $B$ has the higher valence-band edge.  This is suitable for testing an interfacial carrier-transfer scenario in which excited carriers can leave $A_c$ and populate $B_c$, followed by relaxation into $B_v$.

\section{SOC-Derived Diagonal Splitting}

The code does not expose independent spin-splitting parameters.  Instead, the diagonal SOC splitting is derived from material SOC strengths:
\begin{equation}
\lambda_{\mathrm{soc},A},
\qquad
\lambda_{\mathrm{soc},B}.
\end{equation}
The convention used by the code is
\begin{equation}
E_{\uparrow}=E_0+\frac{\Delta}{2},
\qquad
E_{\downarrow}=E_0-\frac{\Delta}{2}.
\end{equation}
For the default qualitative scenario, spin-down CBM lies above spin-up CBM.  Therefore the internal sign convention is
\begin{equation}
\Delta_A = -|\lambda_{\mathrm{soc},A}|,
\qquad
\Delta_B = -|\lambda_{\mathrm{soc},B}|.
\end{equation}
With this sign choice,
\begin{equation}
E_{X,\downarrow}=E_0+\frac{|\lambda_{\mathrm{soc},X}|}{2},
\qquad
E_{X,\uparrow}=E_0-\frac{|\lambda_{\mathrm{soc},X}|}{2}.
\end{equation}
This makes the spin-down band edge higher than the spin-up band edge for both $A$ and $B$, which is the qualitative ordering assumed in the default scenario.

\section{Hamiltonian Terms}

\subsection{Band Hamiltonian}

The diagonal band Hamiltonian is
\begin{equation}
H_{\mathrm{bands}} = \sum_i E_i \ket{i}\bra{i},
\end{equation}
where $\ket{i}$ is a diabatic basis state such as $\ket{A,c,n,\uparrow}$ or $\ket{B,v,m,\downarrow}$.  The energies $E_i$ include the no-SOC band positions and the SOC-derived diagonal splitting.

\subsection{Material-Internal SOC Mixing}

Diagonal SOC splitting alone does not transfer population between spin channels.  To realize the step
\begin{equation}
A_c\uparrow \rightarrow A_c\downarrow,
\end{equation}
the model includes off-diagonal material-internal SOC mixing.  In the simplest same-band-index form,
\begin{equation}
H_{\mathrm{SOC,mix}} =
\sum_{X=A,B}\sum_{b=v,c}\sum_m
\Omega_{X,b}
\ket{X,b,m,\downarrow}\bra{X,b,m,\uparrow}
+\mathrm{h.c.}
\end{equation}
The corresponding software parameters are
\begin{center}
\begin{tabular}{ll}
\toprule
Parameter & Meaning \\
\midrule
\texttt{soc\_mix\_cb\_A} & $A$ conduction-band up/down SOC mixing \\
\texttt{soc\_mix\_cb\_B} & $B$ conduction-band up/down SOC mixing \\
\texttt{soc\_mix\_vb\_A} & $A$ valence-band up/down SOC mixing \\
\texttt{soc\_mix\_vb\_B} & $B$ valence-band up/down SOC mixing \\
\bottomrule
\end{tabular}
\end{center}
The conduction-band mixing parameters are especially important because they allow the elementary spin-mixing step
\begin{equation}
A_c\uparrow \rightarrow A_c\downarrow.
\end{equation}
They also allow SOC-mixed $B_c$ states, so occupation accumulated in $B_c\downarrow$ can indirectly affect $B_c\uparrow$ optical availability through mixed eigenstates.

\subsection{Spin-Conserving Interlayer Hopping}

The interlayer Hamiltonian is spin-conserving.  The main term is conduction-conduction hopping,
\begin{equation}
H_{\mathrm{interlayer}}^{cc}
=
\sum_{m,n,s}
T^{cc}_{mn}
\ket{B,c,n,s}\bra{A,c,m,s}
+\mathrm{h.c.}
\end{equation}
The matrix element is modeled as
\begin{equation}
T^{cc}_{mn}
=
 t_{AB}
 F_{\mathrm{band}}(m,n)
 F_E(E_{A,c,m,s}-E_{B,c,n,s}).
\end{equation}
Here $t_{AB}$ is the interlayer hopping scale, $F_{\mathrm{band}}$ is a band-index overlap filter, and $F_E$ is an energy-mismatch filter.  The band-index filter is
\begin{equation}
F_{\mathrm{band}}(m,n)
=\exp\left[-\frac{(d_m-d_n)^2}{2w_{\mathrm{inter}}^2}\right],
\end{equation}
where the distances are measured from the relevant band edge.  The energy filter is
\begin{equation}
F_E(\Delta E)=
\exp\left[-\frac{(\Delta E)^2}{2\sigma_{\mathrm{hyb}}^2}\right].
\end{equation}
The corresponding parameters are \texttt{tAB}, \texttt{interface\_band\_width}, and \texttt{hybrid\_energy\_width}.  An optional valence-valence hopping term can be enabled by setting \texttt{tAB\_vv} to a nonzero value:
\begin{equation}
H_{\mathrm{interlayer}}^{vv}
=
\sum_{m,n,s}
T^{vv}_{mn}
\ket{B,v,n,s}\bra{A,v,m,s}+\mathrm{h.c.}
\end{equation}
By default \texttt{tAB\_vv}=0.

\section{Laser Pulse and Optical Coupling}

The laser drives spin-conserving intramaterial valence-conduction transitions:
\begin{equation}
X_v s \leftrightarrow X_c s,
\qquad
X\in\{A,B\},\quad s\in\{\uparrow,\downarrow\}.
\end{equation}
The laser Hamiltonian is
\begin{equation}
H_{\mathrm{laser}}(t)
=
F(t)
\sum_{X,m,n,s}
 d^X_{mn}
\ket{X,c,n,s}\bra{X,v,m,s}
+\mathrm{h.c.}
\end{equation}

\subsection{Pulse Shape}

In full-carrier mode,
\begin{equation}
F(t)=A_0
\sin^2\left(\frac{\pi t}{T_{\mathrm{pulse}}}\right)
\cos\left(\frac{\omega t}{\hbar}\right),
\qquad
0<t<T_{\mathrm{pulse}}.
\end{equation}
Outside the pulse window, $F(t)=0$.  In RWA mode, the fast carrier is removed and the drive uses only the envelope.

\subsection{Optical Matrix Elements}

The optical matrix element is modeled as
\begin{equation}
d^X_{mn}=d^X_0 F_{\mathrm{res}}(E_c-E_v)F_{\mathrm{edge}}(m,n).
\end{equation}
The resonance factor is
\begin{equation}
F_{\mathrm{res}}=
\exp\left[-\frac{(E_c-E_v-\hbar\omega)^2}{2\sigma_{\mathrm{opt},X}^2}\right].
\end{equation}
The edge-overlap factor uses distances from CBM and VBM.  If $d_c=0$ denotes the CBM and $d_v=0$ denotes the VBM, then
\begin{equation}
F_{\mathrm{edge}}
=
\exp\left[-\frac{(d_c-d_v)^2}{2w_{\mathrm{opt}}^2}\right].
\end{equation}
This definition intentionally couples CBM-like states most strongly to VBM-like states and avoids the earlier unphysical pairing of CBM with the deepest valence states.


\subsection{Optical Matrix-Element Formula and Input-Parameter Mapping}

Combining the resonance and band-edge overlap factors, the optical matrix element used by the laser term can be written as
\begin{equation}
d^X_{vc}
=
d^X_0
\exp\left[
-\frac{(E_c-E_v-\hbar\omega)^2}{2\sigma_{\mathrm{opt},X}^2}
\right]
\exp\left[
-\frac{(d_c-d_v)^2}{2w_{\mathrm{band}}^2}
\right],
\end{equation}
where $X=A$ or $B$.  The correspondence between the symbols in this expression and the software inputs is:
\begin{center}
\begin{tabular}{ll}
\toprule
Symbol & Input parameter or generated quantity \\
\midrule
$X$ & material label, either $A$ or $B$ \\
$d^A_0$ & \texttt{dA0} \\
$d^B_0$ & \texttt{dB0} \\
$\hbar\omega$ & \texttt{omega\_eV} \\
$\sigma_{\mathrm{opt},A}$ & \texttt{optical\_energy\_width\_A}, or fallback \texttt{optical\_energy\_width} \\
$\sigma_{\mathrm{opt},B}$ & \texttt{optical\_energy\_width\_B}, or fallback \texttt{optical\_energy\_width} \\
$w_{\mathrm{band}}$ & \texttt{band\_overlap\_width} \\
$E_c,E_v$ & generated from gaps, offsets, bandwidths, and SOC-derived diagonal splitting \\
$d_c,d_v$ & generated from local band indices relative to CBM and VBM \\
\bottomrule
\end{tabular}
\end{center}

For a conduction state, $d_c$ is the distance from the CBM in band-index space:
\begin{equation}
\mathrm{CBM}\rightarrow d_c=0,\qquad
\mathrm{next\ CB}\rightarrow d_c=1,\qquad
\mathrm{next\ CB}\rightarrow d_c=2,\ldots
\end{equation}
For a valence state, $d_v$ is the distance from the VBM:
\begin{equation}
\mathrm{VBM}\rightarrow d_v=0,\qquad
\mathrm{next\ lower\ VB}\rightarrow d_v=1,\qquad
\mathrm{next\ lower\ VB}\rightarrow d_v=2,\ldots
\end{equation}
Thus \texttt{band\_overlap\_width} controls how strongly VBM-like states couple to CBM-like states and how quickly this coupling falls off for mismatched band-edge distances.

The state energies $E_c$ and $E_v$ are not direct user inputs.  They are generated by the band model from
\begin{equation}
\{\texttt{gap},\texttt{offset},\texttt{bandwidth\_v},\texttt{bandwidth\_c},\lambda_{\mathrm{soc}}\}.
\end{equation}
Consequently, increasing \texttt{dA0} or \texttt{dB0} increases the material-specific optical coupling scale, but it does not by itself guarantee strong excitation.  The transition must also be close enough to resonance with \texttt{omega\_eV} and sufficiently allowed by the band-index overlap factor.

In compact form, the optical coupling is controlled by four distinct ingredients:
\begin{center}
\begin{tabular}{ll}
\toprule
Input & Role \\
\midrule
\texttt{dA0}, \texttt{dB0} & material-dependent optical coupling scale \\
\texttt{omega\_eV} & photon energy \\
\texttt{optical\_energy\_width}, \texttt{optical\_energy\_width\_A/B} & resonance energy window \\
\texttt{band\_overlap\_width} & band-index overlap window \\
\bottomrule
\end{tabular}
\end{center}

\section{Initial State}

The initial state is not a manually filled diabatic valence configuration.  Instead, the software constructs the full static Hamiltonian
\begin{equation}
H_{\mathrm{static}}=H_{\mathrm{bands}}+H_{\mathrm{SOC,mix}}+H_{\mathrm{interlayer}},
\end{equation}
diagonalizes it,
\begin{equation}
H_{\mathrm{static}}\ket{\phi_n}=\varepsilon_n\ket{\phi_n},
\end{equation}
and fills the lowest $2N$ static eigenstates:
\begin{equation}
\rho(0)=\sum_{n=1}^{2N}\ket{\phi_n}\bra{\phi_n}.
\end{equation}
This avoids a spurious quench at $t=0$ and ensures that the initial density matrix already includes static interlayer and SOC hybridization.

\section{Coherent Time Propagation}

During each time step, the Hamiltonian is evaluated at the midpoint:
\begin{equation}
H_{\mathrm{mid}}=H\left(t+\frac{\Delta t}{2}\right).
\end{equation}
The unitary propagator is
\begin{equation}
U(t+\Delta t,t)
=
\exp\left[-\frac{i}{\hbar}H_{\mathrm{mid}}\Delta t\right].
\end{equation}
The density matrix is updated by
\begin{equation}
\rho(t+\Delta t)=U\rho(t)U^\dagger.
\end{equation}
The implementation computes the exponential by diagonalizing the Hermitian midpoint Hamiltonian:
\begin{equation}
H_{\mathrm{mid}}=V\epsilon V^\dagger,
\qquad
U=V\exp\left[-\frac{i\epsilon\Delta t}{\hbar}\right]V^\dagger.
\end{equation}
For the small matrices used in this toy model, direct dense diagonalization is transparent and sufficiently fast.

\section{Post-Pulse Incoherent Rate Processes}

After the pulse, two phenomenological rate mechanisms are applied.

\subsection{Spin-Conserving Interlayer Downhill Transfer}

The software constructs interlayer transfer channels from the spin-conserving interlayer Hamiltonian.  A channel $i\rightarrow j$ is allowed if
\begin{equation}
E_i>E_j
\end{equation}
and
\begin{equation}
\bra{j}H_{\mathrm{interlayer}}\ket{i}\neq 0.
\end{equation}
The implemented channel rate is
\begin{equation}
\Gamma_{i\rightarrow j}^{\mathrm{inter}}
=
W_{\mathrm{downhill}}
\frac{|\bra{j}H_{\mathrm{interlayer}}\ket{i}|^2}{t_{AB}^2}.
\end{equation}
This describes bath-assisted downhill transfer through spin-conserving interlayer hopping.  It is not an interface SOC spin-flip term.

The normalization by $t_{AB}^2$ is a deliberate toy-model choice.  The coherent Hamiltonian already contains $t_{AB}$ and therefore $t_{AB}$ controls static hybridization and coherent interlayer mixing.  The phenomenological coefficient $W_{\mathrm{downhill}}$ then controls the overall irreversible, bath-assisted post-pulse transfer time scale.  The factor
\begin{equation}
\frac{|\bra{j}H_{\mathrm{interlayer}}\ket{i}|^2}{t_{AB}^2}
\end{equation}
mainly provides channel selectivity from band-index and energy-mismatch filters.  This separation keeps coherent interlayer hybridization and irreversible downhill transfer independently tunable.  A more microscopic Fermi-golden-rule model could instead use $\Gamma\propto |\bra{j}H_{\mathrm{interlayer}}\ket{i}|^2S(\Delta E)$ with an explicit bath spectral function $S$, but such a bath is not part of the present code.  In the present implementation, \texttt{W\_downhill} should therefore be interpreted as an effective phenomenological rate scale rather than a quantity derived uniquely from \texttt{tAB}.

\subsection{Intramaterial $c\rightarrow v$ Relaxation}

Intramaterial relaxation channels have the form
\begin{equation}
X_c s \rightarrow X_v s.
\end{equation}
The channel rate is
\begin{equation}
\Gamma_{c\rightarrow v}^{X}
=
W_{\mathrm{intra},X}
M_{vc}^2
J(E_c-E_v).
\end{equation}
Here $M_{vc}$ reuses the optical matrix-element hierarchy, so different conduction-valence pairs relax at different rates.  The bath-like spectral weight is
\begin{equation}
J(\Delta E)=\frac{\Delta E}{E_{\mathrm{scale}}}\exp\left(-\frac{\Delta E}{E_{\mathrm{scale}}}\right),
\end{equation}
where $E_{\mathrm{scale}}$ is estimated from the material gap and bandwidths.

\subsection{Rate Update Rule}

All one-electron rate channels use an exponential probability
\begin{equation}
p=1-\exp(-\Gamma\Delta t).
\end{equation}
The transferred occupation is
\begin{equation}
\Delta n = p\,n_i(1-n_j).
\end{equation}
The factor $n_i(1-n_j)$ implements Pauli blocking.  This update keeps occupations bounded between 0 and 1 and eliminates the need for a user-facing collision-fraction stability parameter.

\section{Observables and Output Files}

The software outputs several time-dependent observables.  The most important are:
\begin{itemize}[leftmargin=2em]
  \item material- and spin-projected occupations: $N_{A\uparrow}$, $N_{A\downarrow}$, $N_{B\uparrow}$, $N_{B\downarrow}$;
  \item pathway-resolved occupations: $A_v\uparrow$, $A_c\uparrow$, $A_c\downarrow$, $B_c\downarrow$, $B_v\downarrow$, and $B_c\uparrow$;
  \item total spin projections $N_\uparrow$ and $N_\downarrow$;
  \item a magnetic-moment proxy $M=N_\uparrow-N_\downarrow$;
  \item valence and conduction occupation sums;
  \item static eigenbasis low/high occupations;
  \item state-resolved occupation changes relative to pulse end.
\end{itemize}

For output prefix \texttt{interface\_chain\_model}, the standard files are
\begin{lstlisting}
interface_chain_model_projected_material_spin.png
interface_chain_model_pathway_projected_occupations.png
interface_chain_model_projected_spin.png
interface_chain_model_projected_magnetic_moment.png
interface_chain_model_projection_vs_eigen_occupation.png
interface_chain_model_relaxation_diagnostics.png
interface_chain_model_observables.csv
interface_chain_model_state_occupation_change_from_pulse_end_levels.png
interface_chain_model_state_occupation_change_from_pulse_end_levels.csv
\end{lstlisting}
The pathway-resolved plot is the most direct diagnostic for whether the projected occupations support the intended emergent chain-like interpretation.


\section{Matrix-Element Filter Parameters}

Several parameters are not direct observables.  They define smooth model filters for otherwise unknown matrix-element structure.

\begin{itemize}[leftmargin=2em]
  \item \texttt{interface\_band\_width} controls which A/B band indices are connected by spin-conserving interlayer hopping.  It is a width in band-index space, not an energy.
  \item \texttt{hybrid\_energy\_width} controls how strongly interlayer hopping is suppressed when two A/B states are energetically mismatched.
  \item \texttt{band\_overlap\_width} controls which intramaterial valence and conduction states have strong optical and relaxation matrix elements.  It compares distances from VBM and CBM.
  \item \texttt{optical\_energy\_width} controls the tolerance of optical transitions to detuning from the photon energy $\hbar\omega$.
\end{itemize}

Thus \texttt{interface\_band\_width} and \texttt{band\_overlap\_width} are both band-index widths, but they act on different terms: the former acts on interlayer hopping, while the latter acts on intramaterial optical/relaxation matrix elements.  Similarly, \texttt{hybrid\_energy\_width} and \texttt{optical\_energy\_width} are both energy widths, but the former applies to interlayer hybridization and the latter to optical resonance.  The parameters \texttt{dA0} and \texttt{dB0} are not widths; they are material-dependent optical coupling amplitudes multiplying the two optical filters.

\section{Parameter Reference}

\subsection{Band Parameters}

\begin{longtable}{p{0.30\textwidth}p{0.62\textwidth}}
\toprule
Parameter & Meaning \\
\midrule
\texttt{N} & Number of states per material and spin block.  Must be even. \\
\texttt{gap\_A}, \texttt{gap\_B} & No-SOC band gaps of materials $A$ and $B$. \\
\texttt{offset\_A}, \texttt{offset\_B} & Rigid energy shifts of the no-SOC spectra. \\
\texttt{bandwidth\_v\_A}, \texttt{bandwidth\_v\_B} & Valence bandwidths. \\
\texttt{bandwidth\_c\_A}, \texttt{bandwidth\_c\_B} & Conduction bandwidths. \\
\bottomrule
\end{longtable}

\subsection{SOC Parameters}

\begin{longtable}{p{0.30\textwidth}p{0.62\textwidth}}
\toprule
Parameter & Meaning \\
\midrule
\texttt{lambda\_soc\_A}, \texttt{lambda\_soc\_B} & Material-internal SOC scales used to derive diagonal spin splitting. \\
\texttt{soc\_mix\_cb\_A}, \texttt{soc\_mix\_cb\_B} & Off-diagonal up/down SOC mixing in conduction bands. \\
\texttt{soc\_mix\_vb\_A}, \texttt{soc\_mix\_vb\_B} & Off-diagonal up/down SOC mixing in valence bands. \\
\bottomrule
\end{longtable}

\subsection{Interlayer Parameters}

\begin{longtable}{p{0.30\textwidth}p{0.62\textwidth}}
\toprule
Parameter & Meaning \\
\midrule
\texttt{tAB} & Spin-conserving $A_c s\leftrightarrow B_c s$ hopping scale. \\
\texttt{tAB\_vv} & Optional spin-conserving $A_v s\leftrightarrow B_v s$ hopping scale. \\
\texttt{interface\_band\_width} & Band-index overlap width for interlayer hopping. \\
\texttt{hybrid\_energy\_width} & Energy mismatch width for interlayer hopping. \\
\bottomrule
\end{longtable}

\subsection{Laser Parameters}

\begin{longtable}{p{0.30\textwidth}p{0.62\textwidth}}
\toprule
Parameter & Meaning \\
\midrule
\texttt{A0} & Laser amplitude in model units. \\
\texttt{omega\_eV} & Photon energy $\hbar\omega$ in eV. \\
\texttt{pulse\_duration} & Pulse duration in fs. \\
\texttt{carrier} & Either \texttt{full} or \texttt{rwa}. \\
\texttt{dA0}, \texttt{dB0} & Material-dependent optical coupling scales $d^A_0$ and $d^B_0$ in the optical matrix element. \\
\texttt{optical\_energy\_width} & Fallback optical resonance width $\sigma_{\mathrm{opt}}$ used when material-specific widths are not supplied. \\
\texttt{optical\_energy\_width\_A}, \texttt{optical\_energy\_width\_B} & Material-specific optical widths. Blank or \texttt{None} means use the fallback. \\
\texttt{band\_overlap\_width} & Band-edge overlap width $w_{\mathrm{band}}$ for optical and intramaterial relaxation matrix elements. \\
\bottomrule
\end{longtable}

\subsection{Rate and Time Parameters}

\begin{longtable}{p{0.30\textwidth}p{0.62\textwidth}}
\toprule
Parameter & Meaning \\
\midrule
\texttt{W\_downhill} & Effective bath-assisted post-pulse spin-conserving interlayer downhill transfer time scale.  It is kept separate from \texttt{tAB}; \texttt{tAB} enters the coherent Hamiltonian, while \texttt{W\_downhill} controls irreversible transfer. \\
\texttt{W\_intra} & Default intramaterial $c\rightarrow v$ relaxation rate scale. \\
\texttt{W\_intra\_A}, \texttt{W\_intra\_B} & Optional material-specific $c\rightarrow v$ relaxation rates. \\
\texttt{t\_final} & Final simulation time in fs. \\
\texttt{dt} & Time step in fs. \\
\texttt{compare\_time\_ref} & Reference time for state-resolved difference plots.  Default is pulse end. \\
\texttt{compare\_time\_1}, \texttt{compare\_time\_2} & Left and right comparison times for state-resolved level plots. \\
\bottomrule
\end{longtable}

\subsection{Plotting Parameters}

\begin{longtable}{p{0.30\textwidth}p{0.62\textwidth}}
\toprule
Parameter & Meaning \\
\midrule
\texttt{plot\_mode} & Plot either changes (\texttt{delta}) or absolute occupations. \\
\texttt{max\_plot\_points} & Maximum number of plotted time points.  Does not affect propagation. \\
\texttt{delta\_color\_scale} & Manual state-resolved color scale for $\Delta n$.  If $\le 0$, an automatic scale is used. \\
\texttt{delta\_color\_norm} & Either \texttt{linear} or \texttt{symlog}. \\
\texttt{delta\_color\_percentile} & Percentile used for automatic color scaling. \\
\texttt{delta\_colormap} & Diverging colormap for state-resolved occupation changes. \\
\texttt{level\_half\_length}, \texttt{level\_x\_jitter}, \texttt{level\_linewidth}, \texttt{level\_alpha} & Appearance controls for state-resolved energy-level plots. \\
\bottomrule
\end{longtable}

\section{Command-Line Usage}

From the repository root, run without installation by setting \texttt{PYTHONPATH}:
\begin{lstlisting}
PYTHONPATH=. python -m ABinterface --help
\end{lstlisting}
A typical simulation is
\begin{lstlisting}
PYTHONPATH=. python -m ABinterface \
  --N 12 \
  --pulse-duration 20 \
  --t-final 70 \
  --dt 0.02 \
  --carrier full \
  --omega-eV 2.0 \
  --lambda-soc-A 0.05 \
  --lambda-soc-B 0.05 \
  --soc-mix-cb-A 0.02 \
  --soc-mix-cb-B 0.02 \
  --tAB 0.08 \
  --W-downhill 0.01 \
  --W-intra 0.03 \
  --compare-time-1 22 \
  --compare-time-2 50 \
  --delta-color-scale 0.08 \
  --delta-color-norm linear
\end{lstlisting}
After installation,
\begin{lstlisting}
pip install -e .
ABinterface-run --help
\end{lstlisting}

\section{Desktop GUI Usage}

The desktop GUI is based on Tkinter and can be launched by
\begin{lstlisting}
PYTHONPATH=. python -m ABinterface.gui
\end{lstlisting}
or, after installation,
\begin{lstlisting}
ABinterface-gui
\end{lstlisting}
On some Linux systems, the Tkinter bindings must be installed separately:
\begin{lstlisting}
sudo apt install python3-tk
\end{lstlisting}

The GUI uses a tabbed layout:
\begin{center}
\begin{tabular}{ll}
\toprule
Tab & Contents \\
\midrule
Bands & Gaps, offsets, bandwidths, and basis size. \\
SOC + Interlayer & Material SOC splitting, SOC mixing, and spin-conserving interlayer hopping. \\
Laser & Pulse, photon energy, optical matrix elements, and optical widths. \\
Rates + Time & Post-pulse rates and simulation time settings. \\
Plotting & Color scales and state-level plot appearance. \\
\bottomrule
\end{tabular}
\end{center}
The top bar contains the output directory, \texttt{Run simulation}, \texttt{Reset}, and \texttt{Open output}.  The \texttt{Reset} button restores all parameter widgets to \texttt{ModelConfig} defaults.

\section{How to Interpret Results}

A run in which the intended emergent chain-like interpretation is dominant may show the following qualitative behavior:
\begin{enumerate}[leftmargin=2em]
  \item During the pulse, the laser produces $A_v\rightarrow A_c$ and $B_v\rightarrow B_c$ excitations.
  \item After the pulse, material-internal SOC mixing can redistribute projected occupation between $A_c\uparrow$ and $A_c\downarrow$.
  \item Spin-conserving interlayer hopping and the post-pulse downhill channel can move projected occupation between $A_c\downarrow$ and $B_c\downarrow$ when energetically allowed.
  \item Intramaterial $c\rightarrow v$ relaxation can increase $B_v\downarrow$ occupation when $B_c\downarrow$ is populated and final-state phase space is available.
  \item Occupation buildup in $B_c\downarrow$ can influence other $B_c$ channels when $B_c\uparrow$ and $B_c\downarrow$ are SOC-mixed.
\end{enumerate}
The file \texttt{*\_pathway\_projected\_occupations.png} is the most direct way to test this.  The state-resolved level plot is useful for checking whether only a few band-edge states change or whether redistribution extends across multiple bands.

\section{Model Limitations}

This is a toy model.  It is not a first-principles TDDFT, GW, or Boltzmann transport code.  Important limitations include:
\begin{itemize}[leftmargin=2em]
  \item no explicit $k$-resolved band structure;
  \item no microscopic electron-phonon or electron-electron matrix elements;
  \item no self-consistent charge redistribution or band renormalization;
  \item diagonal SOC splitting and SOC mixing are simplified model parameters;
  \item the magnetic moment is a projected spin-population proxy, not a full magnetic observable including orbital contributions.
\end{itemize}
The model is best used to reason about mechanisms, parameter sensitivity, and qualitative dynamical pathways.

\section{Recommended Development Directions}

Potential future extensions include:
\begin{enumerate}[leftmargin=2em]
  \item adding $k$-resolved dispersions and $k$-dependent optical matrix elements;
  \item deriving relaxation rates from explicit phonon or Coulomb spectral functions;
  \item adding optional self-consistent energy shifts from accumulated charge;
  \item separating coherent SOC precession from dephasing-controlled spin relaxation;
  \item adding pathway-resolved current diagnostics between $A_c\downarrow$ and $B_c\downarrow$.
\end{enumerate}

\end{document}
